\subsection{Mecanismo de Recuperación en Dos Etapas y Verificación Bimodal}
\label{subsec:recuperacion_dos_etapas}

La búsqueda sobre grandes colecciones de vídeo enfrenta una disyuntiva clásica entre latencia y precisión: la comparación exhaustiva mediante modelos de visión-lenguaje de alta capacidad resulta computacionalmente inviable en tiempo real, mientras que la recuperación densa basada exclusivamente en similitud del coseno puede verse afectada por coincidencias superficiales y falsos positivos \cite{video_rag2024, eval_mrag2025}. Para resolver este dilema, el sistema adopta una arquitectura en dos etapas jerárquicas: recuperación aproximada de alta cobertura (Etapa 1) y verificación bimodal estructurada (Etapa 2).

\subsubsection{Etapa 1: Recuperación Densa y Deduplicación Temporal}
La consulta multimodal $\mathcal{Q}$ emitida por el usuario es vectorizada en un embedding unitario $\mathbf{q}^{(d)} \in \mathbb{S}^{d-1}$ a la resolución dimensional seleccionada $d$. A continuación, el motor de búsqueda vectorial local \textit{Qdrant} ejecuta una búsqueda de vecinos más cercanos basada en el índice jerárquico de grafos navegables de mundo pequeño (HNSW) con función de distancia coseno:
\begin{equation}
    \mathcal{C}_{\text{raw}} = \text{Top-}k_1 \left( \{ (C_i, S_d(\mathbf{q}, \mathbf{v}_i)) \}_{i=1}^N \right)
\end{equation}
donde el factor de sobremuestreo se establece dinámicamente en $k_1 = \min(5 \cdot k_{\text{target}}, 20)$, garantizando un amplio espacio de candidatos para la etapa posterior.

Para evitar redundancias causadas por cortes visuales muy próximos dentro de una misma escena continua, la lista preliminar $\mathcal{C}_{\text{raw}}$ es sometida a un algoritmo determinista de deduplicación temporal:
\begin{itemize}
    \item \textbf{Fusión de micro-planos adyacentes:} Si dos candidatos recuperados pertenecen al mismo archivo de vídeo y su distancia temporal es inferior a $1,0\text{ s}$, se unifican en un único intervalo envolvente $[t_{\min}, t_{\max}]$.
    \item \textbf{Deduplicación de locuciones repetidas:} Si fragmentos adyacentes comparten una transcripción idéntica o solapada superior al 80\% de coincidencia léxica, se preserva únicamente el candidato con mayor similitud en la primera etapa.
\end{itemize}

El resultado es un conjunto curado $\mathcal{C}_{\text{dedup}}$ de candidatos únicos, ordenados descendentemente por su puntuación vectorial $S_d(\mathbf{q}, \mathbf{v}_i)$.

\subsubsection{Etapa 2: Árbitro Bimodal Multimodal con Gemini 3.5 Flash Lite}
Los mejores candidatos de $\mathcal{C}_{\text{dedup}}$ son transferidos a la segunda etapa de verificación. Para cada candidato $C_i$, el sistema recupera de la caché local en disco los fotogramas clave en formato JPEG y la transcripción de audio alineada $\tau_i$.

El modelo de lenguaje multimodal \textit{Gemini 3.5 Flash Lite} actúa como juez bimodal independiente. Se formula un \textit{prompt} estructurado que incluye:
\begin{enumerate}
    \item La consulta original del usuario (incluyendo imágenes de referencia o transcripción de audio de la consulta).
    \item La evidencia visual del candidato (fotogramas clave secuenciales del fragmento).
    \item La evidencia auditiva textualizada (transcripción de las palabras efectivamente pronunciadas en ese intervalo).
    \item Instrucciones de juicio disjunto e inclusivo: un fragmento se considera relevante si la respuesta a la consulta se encuentra presente en la imagen, en la locución hablada, o en la combinación simultánea de ambas modalidades.
\end{enumerate}

El modelo responde forzando una gramática JSON estricta para cada candidato:
\begin{equation}
    \text{Evaluación}(C_i) = \left\{ \text{relevance\_score} \in [0,0, 1,0], \,\, \text{is\_relevant} \in \{\text{true}, \text{false}\}, \,\, \text{reasoning} \in \text{String} \right\}
\end{equation}

\subsubsection{Discriminación de Falsos Positivos y Ordenación Definitiva}
La puntuación final asignada al fragmento combina la evidencia de ambas etapas y aplica una función de corte estricta:
\begin{equation}
    \text{Score}_{\text{final}}(C_i) = \begin{cases}
        \text{relevance\_score}_i, & \text{si } \text{is\_relevant}_i = \text{true} \land \text{relevance\_score}_i \ge \theta_{\text{rel}} \\
        0,0, & \text{en caso contrario}
    \end{cases}
\end{equation}
donde $\theta_{\text{rel}} = 0,5$. Los candidatos que no superan el criterio del árbitro bimodal son descartados inmediatamente, lo que garantiza que consultas con conceptos inexistentes o trampas visuales no superen la fase de validación (inmunidad total frente a falsos positivos). La lista final de resultados devuelta al usuario se trunca a los $k_{\text{target}}$ mejores elementos, acompañados del razonamiento explícito generado por el modelo.
